\begin{explanationbox}
	\textbf{\textcolor{CExplanation}{一句话直觉}}

	在圆中，垂直于某弦的直径就像一条对称轴，把弦及其所对的弧完全对折重合。

	\medskip
	\textbf{\textcolor{CExplanation}{核心拆解}}
	\begin{description}[style=nextline, leftmargin=2.8em, labelsep=0.5em, itemsep=0.45em]
		\item[\textbf{要点一（垂直条件）：}] 直径与弦互相垂直，这是定理成立的根本前提。
		\item[\textbf{要点二（平分结论）：}] 垂直带来弦被平分，进而弧也被平分。
	\end{description}

	\medskip
	\textbf{\textcolor{CExplanation}{几何本质（圆的对称性）}}

	圆是轴对称图形，任何直径都是圆的对称轴。当直径垂直于弦时，该直径恰好是弦的对称轴，因此弦和弧都被对称地平分。

	\medskip
	\textbf{\textcolor{CExplanation}{代数意义（勾股关系）}}

	在圆心、弦中点和垂足构成的直角三角形中，半径 $R$、半弦长 $\frac{l}{2}$、弦心距 $d$ 满足 $(\frac{l}{2})^2 + d^2 = R^2$，它是圆中用代数手段求值的基础。

	\medskip
	\textbf{\textcolor{CExplanation}{考点价值（高频应用）}}
	\begin{itemize}[leftmargin=*, itemsep=0.5em, topsep=0.4em]
		\item \textbf{考法一（求弦长）：} 已知半径和弦心距，直接代入公式求弦长。
		\item \textbf{考法二（求半径或弦心距）：} 已知弦长和其中一个量，通过勾股定理求另一个。
	\end{itemize}

	\medskip
	\textbf{\textcolor{CExplanation}{顿悟点}}

	\textbf{理解“半径 $R$、半弦长、弦心距 $d$ 构成直角三角形”这一核心结构，几乎所有垂径定理的计算题都围绕它展开。}

	\medskip
	\textbf{\textcolor{CExplanation}{使用场景}}
	\begin{itemize}[leftmargin=*, itemsep=0.5em, topsep=0.4em]
		\item \textbf{场景一（求弦长）：} 在圆中已知圆心到弦的距离和圆半径，求弦长。
		\item \textbf{场景二（求半径或圆心距）：} 已知弦长和另一个条件，求圆的半径或弦心距。
	\end{itemize}
\end{explanationbox}
